\clearpage\chapter{Braid theory}
\section*{1. Overview}
\subsection*{Intuitive}
\paragraph{The problem}

When several strands move from one row of points to another, their over-and-under crossings remember more than the final permutation of the endpoints. Braid theory turns these motions into algebra. Braids can be stacked, undone, represented by generators and relations, closed to make knots, and studied through configuration spaces and mapping classes.



\subsection*{Formal}
The braid group $B_n$ is generated by $\sigma_1,\ldots,\sigma_{n-1}$, where $\sigma_i$ crosses strand $i$ over strand $i+1$. Its defining relations are $\sigma_i\sigma_j=\sigma_j\sigma_i$ when $|i-j|\ge2$ and $\sigma_i\sigma_{i+1}\sigma_i=\sigma_{i+1}\sigma_i\sigma_{i+1}$. Multiplication is stacking braids, and the map $B_n\to S_n$ records the permutation of endpoints.
\section*{2. History}
\subsection*{Historical development}
\begin{historylist}
\paragraph{History and the basic picture}
\bookfigure{knot_braid}{A braid records the ordered crossings of moving strands; closing it produces a knot or link.}

Braid groups were introduced explicitly by Emil Artin in 1925, although Adolf Hurwitz had already encountered their configuration-space meaning in work on monodromy. Artin's 1947 presentation made the groups accessible through generators and relations. Joan Birman's 1974 book connected braids, links, and mapping class groups and became a central reference.

Take $n$ points in a row and connect them to another row. Each strand moves monotonically from left to right. Two diagrams represent the same braid when one can be continuously deformed into the other without strands passing through one another. Over-crossing and under-crossing are part of the data. Stack one braid after another to compose them. Parallel strands give the identity, and reflecting a braid reverses it to give the inverse \cite{braid_theory_ref1,braid_theory_ref4}.

\end{historylist}
\section*{3. Theory}
\subsection*{Core theory}
\paragraph{The braid group and its configuration-space meaning}

The braid group $B_n$ consists of isotopy classes of braids on $n$ strands. Equivalently, let $X$ be a connected manifold of dimension at least two. The unordered configuration space of $n$ distinct points is
\[
 \operatorname{Conf}_n(X)=
 \bigl(X^n\setminus\{x_i=x_j\text{ for some }i\ne j\}\bigr)/S_n.
\]
Its fundamental group is the braid group of $X$. A loop in this configuration space records points moving around one another and returning as an unordered set. If the labels must return to their original positions, the loop lies in the pure braid group.

This interpretation explains why braids compose and why homotopy, rather than a particular drawing, is the correct equivalence relation. For the plane, it gives the classical Artin braid group. In some cases the higher homotopy groups of the configuration space are trivial, making the braid group the main topological information.

\paragraph{Generators and relations}

Let $\sigma_i$ be the braid in which strand $i$ crosses over strand $i+1$, with all other strands straight. The generators $\sigma_1,\ldots,\sigma_{n-1}$ satisfy
\[
 \sigma_i\sigma_j=\sigma_j\sigma_i\quad\text{when }|i-j|\geq2,
\]
because distant crossings do not interact, and
\[
 \sigma_i\sigma_{i+1}\sigma_i=
 \sigma_{i+1}\sigma_i\sigma_{i+1},
\]
because the two three-crossing diagrams are isotopic. The group presentation is
\[
 B_n=\langle\sigma_1,\ldots,\sigma_{n-1}\mid
 \sigma_i\sigma_{i+1}\sigma_i=\sigma_{i+1}\sigma_i\sigma_{i+1},
 \ \sigma_i\sigma_j=\sigma_j\sigma_i\ (|i-j|\geq2)\rangle.
\]
The cubic relation is the braid relation appearing in the Yang--Baxter equation and in Artin groups.

\paragraph{Elementary properties}

$B_1$ is trivial and $B_2\cong\mathbb Z$. The group $B_3$ is infinite and nonabelian and is closely related to the trefoil knot group. Each $B_n$ embeds in $B_{n+1}$ by adding a straight strand, and their increasing union is the infinite braid group $B_\infty$.

Braid groups are torsion-free. For $n\geq3$, they contain a free group on two generators and admit a left-invariant Dehornoy order. Sending every $\sigma_i$ to 1 gives a homomorphism $B_n\to\mathbb Z$ and identifies the abelianization.

\paragraph{Symmetric group and pure braids}

Forget the over-under information but remember where each strand ends. This gives a surjective homomorphism
\[
 B_n\longrightarrow S_n,\qquad \sigma_i\mapsto(i,i+1).
\]
The kernel is the pure braid group $P_n$, in which each strand returns to its own starting point. Pure braids are the fundamental group of the ordered configuration space. They fit into split exact sequences
\[
 1\longrightarrow F_{n-1}\longrightarrow P_n\longrightarrow P_{n-1}\longrightarrow1,
\]
so they can be built as iterated semidirect products of free groups.

\paragraph{Three strands and the modular group}

The group $B_3$ has a distinguished central element. Modding out by its center gives the modular group $\operatorname{PSL}(2,\mathbb Z)$. Thus $B_3$ is a universal central extension of the modular group in the relevant sense. With
\[
 a=\sigma_1\sigma_2\sigma_1,\qquad b=\sigma_1\sigma_2,
\]
the braid relation yields $a^2=b^3$, and the common value is central. This relationship connects braids to modular forms, mapping classes of the punctured torus, and low-dimensional topology \cite{braid_theory_ref2,braid_theory_ref5}.

\paragraph{Closed braids and knots}

Join corresponding top and bottom endpoints of a braid. The result is a link in three-dimensional space, possibly with one component (a knot) or several components. Alexander's theorem says every link is isotopic to the closure of some braid.

Different braids can have the same closure. Markov's theorem gives the exact moves relating braid representatives of the same link: conjugation and stabilization/destabilization. The Jones polynomial was first defined using braid representations and then shown to depend only on the closed link. The braid index is the smallest number of strands needed to represent a link as a closed braid, and it equals the minimum number of Seifert circles in a projection.

\paragraph{Mapping class groups and classification}

Braids can be viewed as mapping classes of a disk with marked points. A braid is a homeomorphism of the punctured disk, considered up to isotopy while fixing the boundary. This makes braid theory part of the study of surfaces.

The Nielsen--Thurston classification divides mapping classes into periodic, reducible, and pseudo-Anosov types. Pseudo-Anosov behavior stretches and folds curves at a definite rate; that rate gives a lower bound for topological entropy. The classification is useful in fluid mixing, where moving rods create braid patterns in space-time trajectories.

\section*{4. Important theorems and results}
\begin{enumerate}[leftmargin=*,itemsep=0.35em]

\item \textbf{Artin presentation.} The braid group is generated by adjacent crossings subject only to far-commutativity and the three-strand braid relations.

\item \textbf{Alexander's theorem.} Every link in three-space is isotopic to the closure of a braid.

\item \textbf{Markov's theorem.} Two braids have isotopic closures precisely when they are related by conjugation and stabilization or destabilization moves.

\item \textbf{Configuration-space theorem.} The braid group is the fundamental group of the configuration space of distinct unordered points.

\item \textbf{Pure-braid extension.} The kernel of $B_n\to S_n$ is $P_n$, and $P_n$ is built from iterated semidirect products of free groups.

\item \textbf{Nielsen--Thurston classification.} Mapping classes, including braids, are periodic, reducible, or pseudo-Anosov, with distinct dynamical behavior.
\end{enumerate}
\noindent\emph{Source for these results:} \cite{braid_theory_ref1}.

\section*{5. Applications}
Braid theory models the motion of distinct strands and records how their crossings compose. The resulting algebra connects configuration spaces with knots, mixing, monodromy, and operators in mathematical physics.
\begin{enumerate}[leftmargin=*,itemsep=0.35em]
\item \textbf{Knot theory.} Closing a braid represents every link. Braid words provide finite data for knot invariants, Markov moves explain when two words describe the same link, and representations lead to polynomial invariants.
\item \textbf{Fluid mixing.} Physical rods, periodic orbits, or almost-invariant sets trace braids. The Nielsen--Thurston type and its entropy estimate how efficiently the stirring produces stretching and mixing.
\item \textbf{Quantum physics.} In two spatial dimensions, exchanging indistinguishable particles can produce braid-group rather than symmetric-group statistics. Anyons use this possibility, and their braid representations are proposed for fault-tolerant quantum computation.
\item \textbf{Algebraic geometry.} Moving roots of a polynomial around one another produces monodromy, represented by a braid. The resulting braid action records how solutions permute when parameters travel around singular values.
\item \textbf{Yang--Baxter equation.} Operators satisfying the braid relation lead to solvable lattice models, quantum groups, and statistical mechanics. The geometric crossing relation becomes an algebraic equation for an operator.
\end{enumerate}


\theoryconnections{braid_theory}{%
\item A braid records paths of moving points, so configuration spaces turn physical motion into loops and group theory records how those loops compose.
\item The braid relation says that sliding one crossing past another does not change the final motion; this is a topological fact expressed as an algebraic equation.
\item Closing a braid produces a link, which is why knot theory and mapping class groups appear naturally \cite{braid_theory_ref1,braid_theory_ref3}.
}{%
\item The same crossing operations model stirring: repeated braid words describe how fluid filaments are stretched and folded, and their growth can estimate mixing.
\item In quantum information and quantum groups, braid generators act as operators whose relations guarantee that different sequences of pairwise exchanges agree.
\item The closure construction also turns algebraic braid data into knot polynomials and other invariants \cite{braid_theory_ref2,braid_theory_ref4,braid_theory_ref5}.
}

\section*{Glossary}
\begin{description}[leftmargin=*,style=nextline,itemsep=0.3em]

\item[\textbf{Braid group.}] The group $B_n$ generated by elementary crossings $\sigma_1,\ldots,\sigma_{n-1}$ subject to far-commutativity and the braid relation; its elements represent ways $n$ strands can wind around one another.

\item[\textbf{Braid relation.}] The identity $\sigma_i\sigma_{i+1}\sigma_i=\sigma_{i+1}\sigma_i\sigma_{i+1}$, which says that sliding one crossing past another produces the same braid; it is the defining cubic relation of the braid group.

\item[\textbf{Pure braid group.}] The kernel of the natural map from the braid group to the symmetric group, consisting of braids in which each strand returns to its original position.

\item[\textbf{Closure of a braid.}] The link obtained by joining each bottom endpoint of a braid to its corresponding top endpoint; Alexander's theorem guarantees that every link arises this way.

\item[\textbf{Configuration space.}] The space of $n$ distinct unordered points in a manifold; its fundamental group is the braid group, giving braids a topological interpretation as loops of moving points.

\item[\textbf{Markov moves.}] The moves (conjugation and stabilization) that relate two different braid words whose closures represent the same link; they are the equivalence criterion for braid representations of knots.

\item[\textbf{Nielsen--Thurston classification.}] A classification of mapping classes---including braids---as periodic, reducible, or pseudo-Anosov, with distinct dynamical behaviors that measure stretching and folding.

\item[\textbf{Modular group.}] The group $\operatorname{PSL}(2,\mathbb Z)$ of $2\times2$ integer matrices with determinant 1 modulo $\pm I$; isomorphic to $B_3$ modulo its center.

\item[\textbf{Symmetric group.}] The group $S_n$ of all permutations of $n$ elements; the image of the braid group under the map forgetting over-under information.

\item[\textbf{Jones polynomial.}] A knot invariant originally defined using braid representations, distinguishing many pairs of knots.

\item[\textbf{Abelianization.}] The quotient of a group by its commutator subgroup, making all generators commute.

\end{description}

\section*{7. Sample Questions}
\emph{Exercise reference:} \cite{braid_theory_ref1}. The cited edition is the source for the exercise set; consult its Exercises section for the official statement and numbering.
\begin{enumerate}[leftmargin=*,itemsep=0.9em]

\item \textbf{Exercise 1.} In the braid group $B_3$, verify the braid relation by computing the permutation in $S_3$ of both sides of $\sigma_1\sigma_2\sigma_1=\sigma_2\sigma_1\sigma_2$.

\emph{Solution.} Each $\sigma_i$ maps to the transposition $(i,i+1)$. Left side: $(1,2)(2,3)(1,2)$. Computing: $(1,2)(2,3)=(1,2,3)$, then $(1,2,3)(1,2)=(1,3)$. Right side: $(2,3)(1,2)(2,3)$. Computing: $(2,3)(1,2)=(1,3,2)$, then $(1,3,2)(2,3)=(1,3)$. Both sides give the transposition $(1,3)$ in $S_3$, consistent with the braid relation.

\item \textbf{Exercise 2.} Compute the image of the braid $\sigma_1\sigma_2\sigma_1^{-1}$ in the symmetric group $S_3$, and determine whether this braid is in the pure braid group $P_3$.

\emph{Solution.} In $S_3$: $\sigma_1\mapsto(1,2)$, $\sigma_2\mapsto(2,3)$, $\sigma_1^{-1}\mapsto(1,2)$. So the image is $(1,2)(2,3)(1,2)=(1,3)$, which is not the identity. Since $P_3$ is the kernel of $B_3\to S_3$, the braid $\sigma_1\sigma_2\sigma_1^{-1}$ is not in $P_3$.

\item \textbf{Exercise 3.} Determine the abelianization of $B_3$ and describe the corresponding homomorphism $B_3\to\mathbb{Z}$.

\emph{Solution.} The abelianization of $B_n$ is $\mathbb{Z}$ for all $n\ge 2$, obtained by sending each generator $\sigma_i$ to $1$. For $B_3=\langle\sigma_1,\sigma_2\mid\sigma_1\sigma_2\sigma_1=\sigma_2\sigma_1\sigma_2\rangle$, the relation becomes $a+b+a=b+a+b$ (writing additively), i.e.\ $a=b$ after cancelling. So both generators map to the same integer. The abelianization is $\mathbb{Z}$, with $\sigma_1,\sigma_2\mapsto 1$.

\item \textbf{Exercise 4.} Find the closure of the braid $\sigma_1\sigma_2^{-1}\in B_3$ and identify the resulting link.

\emph{Solution.} The braid $\sigma_1\sigma_2^{-1}$ has strand 1 crossing over strand 2, then strand 2 crossing under strand 3. Closing this braid connects each top endpoint to its corresponding bottom endpoint. The resulting closed braid has three components (each strand closes to a separate loop), forming the Hopf link union an unlinked unknot, or more precisely the trivial 3-component link $3_1^3$ if no linking occurs. Checking: the permutation is $(1,2)(2,3)=(1,2,3)$, a 3-cycle, so closure gives one component (a knot). Computing more carefully: $\sigma_1$ crosses strands 1,2 and $\sigma_2^{-1}$ crosses strands 2,3 with opposite sign. The closure of this braid is the trefoil knot.

\item \textbf{Exercise 5.} The center $Z(B_3)$ is infinite cyclic, generated by $(\sigma_1\sigma_2)^3$. Verify that $(\sigma_1\sigma_2)^3$ commutes with $\sigma_1$ in $B_3$.

\emph{Solution.} Let $\Delta=(\sigma_1\sigma_2\sigma_1)^2=(\sigma_1\sigma_2)^3$ be the Garside element. We need $\sigma_1\Delta=\Delta\sigma_1$. By the braid relation $\sigma_1\sigma_2\sigma_1=\sigma_2\sigma_1\sigma_2$, so $\Delta=(\sigma_1\sigma_2\sigma_1)^2=\sigma_1\sigma_2\sigma_1\cdot\sigma_1\sigma_2\sigma_1=\sigma_1\sigma_2\sigma_1^2\sigma_2\sigma_1$. Then $\sigma_1\Delta=\sigma_1^2\sigma_2\sigma_1^2\sigma_2\sigma_1$. Similarly $\Delta\sigma_1=\sigma_1\sigma_2\sigma_1^2\sigma_2\sigma_1^2$. Using the braid relation repeatedly shows both expressions are equal. Indeed $\Delta$ is central in $B_3$, a well-known fact.

\end{enumerate}
\section*{8. References}
\begin{thebibliography}{99}
\bibitem{braid_theory_ref1} J. S. Birman, \emph{Braids, Links, and Mapping Class Groups}, Princeton University Press, 1974.
\bibitem{braid_theory_ref2} C. Kassel and V. Turaev, \emph{Braid Groups}, Springer, 2008.
\bibitem{braid_theory_ref3} A. Hatcher, \emph{Algebraic Topology}, Cambridge University Press, 2002.
\bibitem{braid_theory_ref4} V. F. R. Jones, ``A polynomial invariant for knots via von Neumann algebras,'' \emph{Bulletin of the American Mathematical Society}, Vol.~12, No.~1, 1985.
\bibitem{braid_theory_ref5} C. Kassel, \emph{Quantum Groups}, Springer, 1995.
\end{thebibliography}
